\documentclass[conference]{IEEEtran}
\IEEEoverridecommandlockouts
% The preceding line is only needed to identify funding in the first footnote. If that is unneeded, please comment it out.
\usepackage{cite}
\usepackage{amsmath,amssymb,amsfonts}
\usepackage{algorithmic}
\usepackage{graphicx}
\usepackage{textcomp}
\usepackage{xcolor}
\usepackage{hyperref}
\def\BibTeX{{\rm B\kern-.05em{\sc i\kern-.025em b}\kern-.08em
    T\kern-.1667em\lower.7ex\hbox{E}\kern-.125emX}}
\begin{document}

\title{Laboratório 2 - Unidade Operativa Uniciclo RISC-V32I\\}

\author{\IEEEauthorblockN{1\textsuperscript{st} Artur Henrique do Carmo Rosa}
\IEEEauthorblockA{\textit{Departamento de Ciências da Computação} \\
\textit{Universidade de Brasília}\\
Brasília, Brazil \\
222028738@aluno.unb.br}
\and
\IEEEauthorblockN{2\textsuperscript{nd} Daniel Dezan Baby}
\IEEEauthorblockA{\textit{Departamento de Ciências da Computação} \\
\textit{Universidade de Brasília}\\
Brasília, Brazil \\
232046097@aluno.unb.br}
\and
\IEEEauthorblockN{3\textsuperscript{rd} João Filipe Araújo}
\IEEEauthorblockA{\textit{Departamento de Ciências da Computação} \\
\textit{Universidade de Brasília}\\
Brasília, Brazil \\
231013574@aluno.unb.br}
\and
\IEEEauthorblockN{4\textsuperscript{th} Leonardo Tomé Sampaio}
\IEEEauthorblockA{\textit{Departamento de Ciências da Computação} \\
\textit{Universidade de Brasília}\\
Brasília, Brazil \\
222005448@aluno.unb.br}
\and
\IEEEauthorblockN{5\textsuperscript{th} Luís Eduardo Curi Serra}
\IEEEauthorblockA{\textit{Departamento de Ciências da Computação} \\
\textit{Universidade de Brasília}\\
Brasília, Brazil \\
251033574@aluno.unb.br}
}

\maketitle

\begin{abstract}
% TODO (grupo): completar o abstract com a visão geral de todo o trabalho ao final, quando todos os módulos estiverem prontos.
No presente trabalho é apresentado o desenvolvimento da Unidade Operativa Uniciclo da arquitetura
RISC-V32I, implementada no ambiente Quartus II. O projeto contempla os módulos que compõem o caminho de
dados (\textit{datapath}), as memórias de instrução e de dados em organização Harvard, a unidade de
controle principal e a unidade de controle da ULA, de forma a executar automaticamente um conjunto de
instruções de 32 bits pré-compiladas no formato MIF. O trabalho inclui a descrição da arquitetura
desenvolvida, do conjunto de instruções suportado, dos procedimentos e métodos adotados, dos resultados
obtidos por meio do ambiente de simulação Waveform Editor e da análise de desempenho das unidades
funcionais, seguidos de uma análise crítica dos resultados.
\end{abstract}

\begin{IEEEkeywords}
RISC-V, Uniciclo, FPGA, Quartus, caminho de dados, unidade de controle, multiplexador
\end{IEEEkeywords}

%-----------------------------------------------------------------------

\section{Divisão de tarefas}

\subsection{Objetivos e Introdução (Teoria da Arquitetura)}
Seção ii. Objetivos: resumo sucinto dos objetivos do laboratório.
Seção iii. Introdução (parte 1): visão geral da arquitetura RISC-V32I, filosofia RISC, organização
Uniciclo e caminho de dados.

\subsection{Introdução (Caminho de Dados e Memórias)}
Seção iii. Introdução (parte 2): detalhamento dos módulos do caminho de dados, memórias de instrução e
dados em organização Harvard, banco de registradores, somador, gerador de imediato e ULA.

\subsection{Introdução (Unidade de Controle e Multiplexadores)}
Seção iii. Introdução (parte 3): unidade de controle principal, unidade de controle da ULA, sinais de
controle e multiplexadores do caminho de dados.

\subsection{Procedimentos e Métodos}
Seção iv. Materiais e Métodos: ferramentas utilizadas (Quartus II, Waveform Editor), descrição da
implementação de cada módulo e metodologia de sincronização e integração.

\subsection{Resultados experimentais}
Seção v. Resultados: simulações no Waveform Editor, conteúdo dos registradores, sinais de controle e
análise de desempenho de cada unidade funcional.

\subsection{Análise de resultados, Conclusões e Revisão Geral}
Seção vi. Discussão e Conclusões: análise qualitativa e quantitativa, gargalos identificados, pontos
críticos e conclusão geral do laboratório.

%-----------------------------------------------------------------------

\section{Objetivos}
O objetivo deste laboratório foi desenvolver a Unidade Operativa Uniciclo da arquitetura RISC-V32I,
utilizando o ambiente Quartus II Altera, de forma a montar os módulos principais do caminho de dados
(\textit{datapath}) capazes de processar instruções de 32 bits pré-compiladas no formato MIF, fazendo uso
de memórias de instrução e de dados em organização Harvard. Além disso, o trabalho buscou reforçar os
conceitos estudados na disciplina de Organização e Arquitetura de Computadores, proporcionando uma
compreensão mais prática sobre o funcionamento do caminho de dados, da unidade de controle e da análise de
desempenho de um processador Uniciclo.

\section{Introdução}

\subsection{Arquitetura RISC-V32I Uniciclo}
% TODO

\subsection{Contador de Programa e Memória de Instruções}
% TODO

\subsection{Banco de Registradores}
% TODO

\subsection{Unidade Lógica e Aritmética (ULA) e seu Controle}
% TODO

\subsection{Somador}
% TODO

\subsection{Gerador de Imediato}
% TODO

\subsection{Memória de Dados}
% TODO

\subsection{Unidade de Controle Principal}
A Unidade de Controle Uniciclo é o bloco responsável por observar o campo opcode da instrução, com 7 bits
de largura, e decidir quais sinais devem ser ativados no restante do caminho de dados para que aquela
instrução específica seja executada corretamente \cite{patterson2020computer}. Diferente de um registrador
ou de uma memória, ela não guarda nenhum valor de um ciclo para o outro: assim que a memória de instruções
entrega um novo opcode na sua entrada, a saída já reflete os novos sinais de controle, respeitando apenas
o tempo de propagação da lógica combinacional utilizada na sua implementação.

Os sinais produzidos por esse bloco, junto com o efeito de cada um sobre o caminho de dados, são os
seguintes:

\begin{itemize}
    \item RegWrite: habilita a escrita no banco de registradores;
    \item MemRead: habilita a leitura da memória de dados;
    \item MemWrite: habilita a escrita na memória de dados;
    \item MemtoReg: seleciona, por meio de um multiplexador, se o valor gravado no registrador de destino
    vem do resultado da ULA ou do dado lido da memória;
    \item ALUSrc: seleciona, também por meio de um multiplexador, se a segunda entrada da ULA é o conteúdo
    do segundo registrador lido ou o valor imediato gerado pelo gerador de imediato;
    \item Branch: indica que a instrução é um desvio condicional, sinal que é combinado com a saída Zero
    da ULA para decidir o próximo valor do contador de programa;
    \item ALUOp: indica de forma resumida o tipo de operação que a ULA deve realizar, sendo depois
    combinado com os campos funct3 e funct7 pela unidade de controle da ULA para definir a operação final.
\end{itemize}

A Tabela \ref{tab:controle} apresenta a tabela verdade da unidade de controle, relacionando cada grupo de
instruções, identificado por seu opcode, aos valores dos sinais de controle gerados
\cite{patterson2020computer}. O valor X (don't care) indica situações em que o sinal não influencia o
resultado da instrução, pois seu efeito não é utilizado naquele caso. O campo ALUOp utiliza 2 bits para
identificar a categoria da operação a ser realizada pela ULA: 00 para acessos à memória, 01 para desvios e
10 para operações lógico-aritméticas, sendo o detalhamento da operação final feito pela unidade de
controle da ULA a partir dos campos funct3 e funct7.

\begin{table}[!t]
\caption{Tabela verdade da Unidade de Controle}
\begin{center}
\begin{tabular}{c|c|c|c|c|c|c|c|c}
\hline \hline
\textbf{Instr.} & \textbf{Opcode} & \textbf{ALUSrc} & \textbf{Mem2Reg} & \textbf{RegWrite} & \textbf{MemRead} & \textbf{MemWrite} & \textbf{Branch} & \textbf{ALUOp} \\
\hline
Tipo-R & 0110011 & 0 & 0 & 1 & 0 & 0 & 0 & 10 \\
lw & 0000011 & 1 & 1 & 1 & 1 & 0 & 0 & 00 \\
sw & 0100011 & 1 & X & 0 & 0 & 1 & 0 & 00 \\
beq & 1100011 & 0 & X & 0 & 0 & 0 & 1 & 01 \\
\hline \hline
\end{tabular}
\end{center}
\label{tab:controle}
\end{table}

Esses quatro grupos representam o modelo de referência apresentado em sala \cite{patterson2020computer}. A
mesma estrutura combinacional foi utilizada para reconhecer também os demais opcodes do conjunto de
instruções adotado no projeto, atribuindo a cada um deles a combinação adequada de sinais conforme o
formato correspondente.

\subsection{Multiplexador 2:1}
O multiplexador 2 para 1 é o elemento combinacional mais simples utilizado no caminho de dados: ele recebe
duas entradas de dados e, dependendo do valor de um único bit de seleção, decide qual delas aparece na
saída \cite{patterson2020computer}. A Tabela \ref{tab:mux} resume seu comportamento.

\begin{table}[!t]
\caption{Tabela verdade do Multiplexador 2:1}
\begin{center}
\begin{tabular}{c|c}
\hline \hline
\textbf{Seleção} & \textbf{Saída} \\
\hline
0 & Entrada 0 \\
1 & Entrada 1 \\
\hline \hline
\end{tabular}
\end{center}
\label{tab:mux}
\end{table}

Em vez de montar um circuito seletor diferente para cada ponto do caminho de dados onde havia essa
necessidade, optou-se por um único módulo genérico e parametrizável quanto à largura do barramento, depois
instanciado nos três pontos do caminho de dados Uniciclo em que existe seleção entre duas fontes de dados.
O primeiro deles escolhe entre o dado lido do segundo registrador e o valor imediato, de acordo com o
sinal ALUSrc. O segundo escolhe entre o resultado da ULA e o dado lido da memória, de acordo com o sinal
MemtoReg, definindo o que será escrito de volta no banco de registradores. O terceiro escolhe entre o
valor seguinte do contador de programa e o endereço de desvio calculado, de acordo com a combinação dos
sinais Branch e Zero.

\section{Procedimentos e Métodos}

\subsection{Ferramentas Utilizadas}
% TODO

\subsection{Contador de Programa e Memória de Instruções}
% TODO

\subsection{Banco de Registradores}
% TODO

\subsection{Unidade Lógica e Aritmética e seu Controle}
% TODO

\subsection{Somador}
% TODO

\subsection{Gerador de Imediato}
% TODO

\subsection{Memória de Dados}
% TODO

\subsection{Unidade de Controle Principal}
A unidade de controle principal foi descrita como um bloco de lógica combinacional, utilizando uma
estrutura de decisão sobre o campo opcode de 7 bits extraído diretamente da instrução lida da memória.
Para cada opcode suportado, foi definido um conjunto fixo de valores para os sinais de saída (RegWrite,
MemRead, MemWrite, MemtoReg, ALUSrc, Branch e ALUOp), seguindo a tabela verdade apresentada na
introdução. Para opcodes não previstos, todos os sinais de habilitação foram mantidos desligados por
padrão, evitando que uma instrução inválida cause algum efeito indesejado no restante do caminho de dados.

A validação desse módulo foi feita de forma isolada no Waveform Editor do Quartus, aplicando-se cada
opcode como entrada e conferindo, um a um, se a combinação de sinais gerada na saída correspondia
exatamente à tabela de projeto antes da integração com o restante do caminho de dados.

\subsection{Multiplexador 2:1}
O módulo do multiplexador foi escrito de forma genérica, com a largura do barramento de dados como
parâmetro, o que permitiu reaproveitá-lo nos três pontos de seleção do caminho de dados sem reescrever a
lógica em cada lugar. Essa escolha reduz a duplicação de código e facilita a verificação, já que o
comportamento do multiplexador só precisa ser validado uma vez para que a confiança se estenda a todas as
instâncias seguintes. A validação foi feita no Waveform Editor, alternando o sinal de seleção entre 0 e 1
para diferentes combinações de entrada e conferindo se a entrada correta era propagada para a saída em
cada caso.

\section{Resultados experimentais}

\subsection{Contador de Programa e Memória de Instruções}
% TODO

\subsection{Banco de Registradores}
% TODO

\subsection{Unidade Lógica e Aritmética e seu Controle}
% TODO

\subsection{Somador}
% TODO

\subsection{Gerador de Imediato}
% TODO

\subsection{Memória de Dados}
% TODO

\subsection{Unidade de Controle Principal e Multiplexador 2:1}
A verificação de ambos os módulos foi feita por meio de \textit{testbenches} próprios, escritos em Verilog
e executados no simulador, que comparam automaticamente a saída obtida com o valor esperado em cada caso de
teste e reportam o resultado.

Para a unidade de controle, o \textit{testbench} aplicou, um a um, os opcodes correspondentes a todos os
grupos de instruções suportados pelo projeto, além de um opcode inválido para verificar o comportamento
padrão. A Tabela \ref{tab:result_controle} resume os casos testados, todos aprovados.

\begin{table}[!t]
\caption{Casos de teste da Unidade de Controle}
\begin{center}
\begin{tabular}{c|c|c}
\hline \hline
\textbf{Opcode (bin)} & \textbf{Opcode (hex)} & \textbf{Resultado} \\
\hline
0110011 & 33 & PASS \\
0010011 & 13 & PASS \\
0000011 & 03 & PASS \\
0100011 & 23 & PASS \\
1100011 & 63 & PASS \\
1101111 & 6F & PASS \\
1100111 & 67 & PASS \\
0110111 & 37 & PASS \\
0010111 & 17 & PASS \\
1111111 & 7F & PASS \\
\hline \hline
\end{tabular}
\end{center}
\label{tab:result_controle}
\end{table}

Como a unidade de controle é um circuito puramente combinacional, a forma de onda gerada na simulação
apresenta os sinais de saída constantes para cada opcode aplicado, sem transições intermediárias, o que é
o comportamento esperado para esse tipo de lógica.

Para o multiplexador, o \textit{testbench} aplicou diferentes combinações de dados de 32 bits nas duas
entradas, variando o sinal de seleção entre 0 e 1, e conferiu se a saída reproduzia exatamente a entrada
selecionada em cada caso. Foram usados valores representativos como AAAAAAAA, 55555555, DEADBEEF e
CAFEBABE, e todos os casos foram aprovados, confirmando o funcionamento correto do módulo para qualquer
dado de entrada.

% TODO: inserir capturas de tela do Waveform Editor com as formas de onda da unidade de controle e do
% multiplexador, para complementar as tabelas acima.

\subsection{Análise de Desempenho das Unidades Funcionais}
% TODO (grupo): tabela de tempos de execução por módulo, valor absoluto e percentual, identificação dos
% gargalos e estimativa do período de clock.

\section{Análise de Resultados}
% TODO (grupo): completar com a discussão dos demais módulos e da integração completa.
A unidade de controle principal e o multiplexador atenderam ao comportamento esperado em todos os casos de
teste, gerando os sinais corretos para cada opcode e propagando o dado certo em função do sinal de seleção.
A verificação automática por \textit{testbench} foi importante porque permitiu confirmar o comportamento de
cada módulo de forma isolada, antes da integração com o restante do caminho de dados, reduzindo a chance de
erros difíceis de localizar em uma etapa posterior.

\section{Conclusão}
% TODO (grupo): conclusão geral do laboratório, a ser escrita quando todos os módulos estiverem prontos e
% integrados.
A construção da unidade de controle principal e do multiplexador mostrou, na prática, como a decodificação
do opcode pode ser isolada em um bloco puramente combinacional e como um único componente seletor genérico
pode ser reaproveitado em vários pontos do caminho de dados. Essa organização modular deixou o projeto mais
claro e facilitou a verificação individual de cada parte antes da integração da unidade operativa completa.

\section*{Agradecimentos}



\begin{thebibliography}{00}
\bibitem{referenceGuide}
F. Vidal,
\textit{RISC-V Reference Guide v25}, 2025.

\bibitem{patterson2020computer}
D.~A. Patterson e J.~L. Hennessy,
\textit{Computer Organization and Design RISC-V Edition: The Hardware Software Interface}.
Morgan Kaufmann, 2020.

\end{thebibliography}

\end{document}